% Gerado por scripts/etapa3_comparacao.py — requer \usepackage{booktabs}
\begin{table}[htbp]
\centering
\small
\caption{Combinação HW + XGBoost — métricas de acurácia e de custo}
\label{tab:comb_metricas}
\begin{tabular}{lrrrrrrrr}
\toprule
Quantis & MAE & RMSE & Pinball 80\% & Pinball 90\% & Pinball 95\% & Custo 1:4 & Custo 1:9 & Custo 1:19 \\
\midrule
Empírico & 1,245 & 2,287 & 0,530 & 0,380 & 0,249 & 2,620 & 3,768 & 4,958 \\
Normal & 1,245 & 2,287 & 0,539 & 0,378 & 0,252 & 2,668 & 3,744 & 4,981 \\
\bottomrule
\end{tabular}
\par\smallskip
{\footnotesize Nota: teste de 90 dias (26/01 a 24/04/2016), 5.748 séries. MAE e RMSE avaliam a previsão pontual e são iguais nos dois métodos de quantis. Pinball avalia o quantil previsto; o custo avalia a decisão de estoque. custo por série-dia com custo de sobra = 1; razão 1:4, 1:9 e 1:19 (sobra:falta) corresponde aos níveis de serviço de 80\%, 90\% e 95\%. Estoque = quantil arredondado para o inteiro mais próximo.}
\end{table}
